\section{Inventory System}
\label{sec:inventory-system}

\subsection{Overview}

The Inventory System manages all player-owned resources outside of combat. 
It is split by item type into dedicated sub-inventories, but presented as a single, tabbed UI inside the Hideout.

\begin{itemize}
	\item \textbf{Per item type a dedicated inventory:} each item family (Parts, Remains, Runes, Gear, Modules) has its own capacity and stack behaviour.
	\item \textbf{Single Inventory Menu:} all sub-inventories are accessible via tabs in one consolidated screen for clarity.
	\item \textbf{Hero Gear Slots:} each Hero has 6 equipment slots that are filled from the Gear inventory:
	\begin{itemize}
		\item Head, Chest, Gloves, Leg, Boots, Amulet.
	\end{itemize}
	\item \textbf{Hideout-only Access:} inventory management (equipping, moving, crafting usage) is only available in the Hideout, not during Maps.
\end{itemize}

\textbf{Inventory Categories:}
\begin{itemize}
	\item \emph{Items} (Gear, Modules, Runes): Individual, building-related objects with identity and decision-making quality.
	\item \emph{Materials} (Parts, Remains): Stackable resources that primarily serve as input for crafting and upgrading.
\end{itemize}

The Inventory System itself does not generate items or currencies and does not determine drop chances. 
It acts as the persistent container and organiser for all loot generated by the Loot System and all results created by the Crafting System.

% RISK: Multiple sub-inventories increase UI complexity; needs strong visual hierarchy in the tabbed layout.
% RISK: Hideout-only access can feel restrictive if players want to inspect or adjust gear mid-run.

\subsection{Resource Types}

The game distinguishes several item families and currencies. Each family has its own stack rules and usage pattern.

\subsubsection*{Item Types}

\begin{center}
	\begin{tabular}{p{2cm}p{4cm}p{2cm}p{3.5cm}p{4cm}}
		\textbf{Type} & \textbf{Description} & \textbf{Stack Limit} & \textbf{Primary Source} & \textbf{Primary Use} \\
		\hline
		Parts &
		Leftover body parts and fragments from monsters; main generic crafting material. &
		100 &
		Enemy drops (all Maps) &
		Baseline crafting inputs for gear and other items. \\
		\hline
		Remains &
		Monster drops representing meta-materials in various rarities (e.g. ash, glitter dust, boss cores). &
		999 &
		Enemy drops (rarity scaling with monster/map difficulty) &
		Advanced and meta-level crafting operations; higher-tier recipes and reforging. \\
		\hline
		Runes &
		Craftable endgame items. &
		10 &
		Crafting results, rare drops &
		Socketed or slotted into systems (e.g. endgame build customisation). \\
		\hline
		Gear &
		Equipable items for Heroes (weapons/armour/charms, abstracted into slots). &
		1 (no stacking) &
		Crafting results, rare drops, rewards &
		Direct Hero power via equipment: stats, affixes, unique effects. \\
		\hline
		Modules &
		Craftable/findable items for configuring and augmenting skills. &
		10 &
		Crafting, targeted drops, rewards &
		Modify skills: behaviour, triggers, numbers, interactions. \\
	\end{tabular}
	%	\caption{Item Types and Stack Rules}
	\label{tab:item-types}
\end{center}

\subsubsection*{Currencies}

\begin{itemize}
	\item \textbf{Gold}
	\begin{itemize}
		\item Main currency.
		\item Almost all Hideout activities (crafting, reforging, respecs, Map-device usage, etc.) cost Gold.
		\item Acts as the primary brake against endless crafting loops.
	\end{itemize}
	\item \textbf{Crystals}
	\begin{itemize}
		\item Secondary, rarer currency.
		\item Placeholder for high-end or structural unlocks (e.g. endgame systems, Hero unlocks).
		% TODO: Clarify exact use of Crystals (hero unlocks vs. endgame-only sinks) in a later balance pass.
	\end{itemize}
\end{itemize}

% RISK: If both Gold and Crystals gate similar things, players may experience redundant currencies.
% TODO: Define a clear separation: Gold = operational costs, Crystals = strategic/unlock costs.

\subsection{Inventory Structure \& Rules}

\subsubsection*{Global Inventory and Sub-Inventories}

\begin{itemize}
	\item \textbf{Single player inventory:} currently, each player has exactly one main inventory; there are no separate stash tabs or alternate character inventories.
	\item \textbf{Per-type sub-inventories:} within this main inventory, each item type has its own sub-inventory with:
	\begin{itemize}
		\item its own capacity model (stack limits, number of stacks),
		\item type-specific UI representation (tab, iconography),
		\item type-specific sorting/filtering options (future).
	\end{itemize}
	\item \textbf{Hero Equipment:} 
	\begin{itemize}
		\item Heroes do not have separate bags; they only have 6 equipment slots (Head, Chest, Arm, Leg, Boots, Charm).
		\item Equipping and unequipping gear always transfers items between Hero slots and the global Gear inventory.
	\end{itemize}
\end{itemize}

\subsubsection*{Access and Restrictions}

\begin{itemize}
	\item \textbf{Hideout-only:} 
	\begin{itemize}
		\item Inventory management is only available in the Hideout (Hub).
		\item During Maps, the player cannot:
		\begin{itemize}
			\item move items,
			\item change gear,
			\item initiate crafting/recycling actions.
		\end{itemize}
	\end{itemize}
	\item \textbf{During Maps:}
	\begin{itemize}
		\item Items and currencies are collected, but not manipulated.
		\item Gear decisions and build changes are postponed until return to the Hideout.
	\end{itemize}
\end{itemize}

% RISK: Strong restriction to Hideout can frustrate players who want to inspect or theorycraft mid-run.
% TODO: Consider a read-only inventory view during Maps for inspection without interaction.

\subsection{Loot Pickup \& Interaction}

\subsubsection*{Drops}

\begin{itemize}
	\item \textbf{Enemies drop:}
	\begin{itemize}
		\item Parts,
		\item Remains,
		\item Currencies (primarily Gold).
	\end{itemize}
	\item \textbf{Drop presentation:}
	\begin{itemize}
		\item Loot appears as neutral bags on the ground.
		\item Bags are colour-coded by rarity.
	\end{itemize}
\end{itemize}

\subsubsection*{Pickup Behaviour}

\begin{itemize}
	\item \textbf{Manual Pickup:}
	\begin{itemize}
		\item The player clicks on loot bags on the screen.
		\item Each pickup triggers visual and audio feedback, plus a small text log at the side summarising what was collected.
	\end{itemize}
	\item \textbf{Optional „Pick All“:}
	\begin{itemize}
		\item After each Wave, an optional „Pick All“ function can collect all remaining loot at once.
		% ASSUMPTION (70\%): This is a convenience option and may be gated by research or settings later.
	\end{itemize}
	\item \textbf{Currencies:}
	\begin{itemize}
		\item Currency drops from enemies are collected automatically on pickup (no separate item objects).
	\end{itemize}
\end{itemize}

\subsubsection*{Loot Usefulness \& Filtering}

\begin{itemize}
	\item \textbf{No loot filter:}
	\begin{itemize}
		\item There is no explicit loot filter planned.
		\item Design assumption: \emph{every drop is useful} as input into crafting or progression, even if indirectly.
	\end{itemize}
\end{itemize}

% RISK: Without a loot filter, visual clutter and cognitive load can become problematic at high kill rates.
% TODO: Later consider minimal quality-of-life options (e.g. stacking on ground, combined pickup messages).

\subsection{Economic Role of the Inventory}

The Inventory System itself is intentionally designed as a \emph{comfort container}, not as a primary economic pressure tool:

\begin{itemize}
	\item \textbf{No hard space pressure:}
	\begin{itemize}
		\item The main friction in the economy is not inventory space, but the availability of currencies and materials (especially Gold and high-tier Remains).
		\item Item stack limits are present but relatively generous, mainly to maintain readability and avoid infinite hoarding without cost.
	\end{itemize}
	\item \textbf{Economic pressure via costs, not slots:}
	\begin{itemize}
		\item Gold acts as the main brake on crafting and Hideout actions.
		\item Rare currencies (Crystals, late-game Remains) gate high-impact upgrades and endgame systems.
	\end{itemize}
\end{itemize}

% RISK: If inventory has almost no limitations and all loot is useful, players may hoard everything and rarely make meaningful discard decisions.
% TODO: Decide later whether soft sinks (e.g. bulk-conversion, diminishing returns) are needed to keep inventories manageable without hard slot limits.
